\documentclass[11pt]{article}
\usepackage[T1]{fontenc}
\usepackage[utf8]{inputenc}
\usepackage[margin=1in]{geometry}
\usepackage{amsmath,amssymb}
\usepackage{microtype}
\usepackage[hidelinks]{hyperref}
\newcommand{\Pb}{\mathbb P}
\DeclareMathOperator{\Beta}{Beta}
\DeclareMathOperator{\BR}{BR}
\allowdisplaybreaks[1]
\title{Two-Seller Competition: Model and Extinction Questions}
\author{Progress note --- transcription of handwritten notes}
\date{October 7, 2026}

\begin{document}
\maketitle

This note transcribes the three handwritten pages into English. The initial
intuition and its subsequent qualification are presented in their original
order.

\section{Elements of the model}

\begin{itemize}
\item Two symmetric providers, indexed by $j\in\{A,B\}$.
\item Two products, indexed by $i\in\{1,2\}$.
\item Success probabilities: $0<p_1<p_2<1$.
\item Costs: $c_1<c_2$.
\item Revenue: $R>c_2$.
\end{itemize}

\paragraph{Customer beliefs.}
The customer observes the state
\[
s=(S_A,F_A,S_B,F_B),
\]
where $S_j$ and $F_j$ count provider $j$'s successes and failures. Her
beliefs are
\begin{align}
\theta_A\mid S_A,F_A&\sim\Beta(1+S_A,1+F_A),\\
\theta_B\mid S_B,F_B&\sim\Beta(1+S_B,1+F_B).
\end{align}
For the first model, we assume that \emph{both sellers observe the full
state $s$}.

\paragraph{Customer choice.}
The customer uses Thompson sampling (TS), drawing independently from the
two belief distributions and choosing
\begin{equation}
y_t=
\begin{cases}
A,&\widetilde\theta_A>\widetilde\theta_B,\\
B,&\widetilde\theta_B>\widetilde\theta_A.
\end{cases}
\end{equation}
Thus the selection probabilities are
\begin{equation}
D_A(s)=\Pb(\widetilde\theta_A>\widetilde\theta_B\mid s),
\qquad D_B(s)=1-D_A(s).
\end{equation}
Ties have probability zero under the continuous Beta distributions.

\section{Timing, state transitions, and strategies}

\paragraph{Timing.}
\begin{enumerate}
\item Everyone observes $s$.
\item The customer chooses a seller using Thompson sampling.
\item Only the selected provider chooses product $i$.
\item The customer observes the outcome, but not the product.
\item The selected provider's state is updated.
\end{enumerate}

\paragraph{State transitions.}
Superscripts $s$ and $f$ denote success and failure, respectively:
\begin{equation}
\begin{aligned}
T_A^s(s)&=(S_A+1,F_A,S_B,F_B),\\
T_A^f(s)&=(S_A,F_A+1,S_B,F_B),\\
T_B^s(s)&=(S_A,F_A,S_B+1,F_B),\\
T_B^f(s)&=(S_A,F_A,S_B,F_B+1).
\end{aligned}
\end{equation}

\paragraph{Strategies.}
The sellers' product choices, conditional on selection, are
\begin{equation}
\pi_A(s)\in\{1,2\},\qquad \pi_B(s)\in\{1,2\}.
\end{equation}
Here $\pi_j$ denotes a product-choice strategy, as in the handwritten notes.

\section{Bellman equation and equilibrium}

\paragraph{Bellman equation.}
Fix seller $B$'s strategy $\pi_B$. Let $V_A(s)$ be seller $A$'s optimal
discounted value before customer selection, and let $\gamma\in(0,1)$
denote the discount factor. Then
\begin{equation}\label{eq:bellman}
\begin{aligned}
V_A(s)={}&D_A(s)\max_{i\in\{1,2\}}\Bigl\{
R-c_i+\gamma\bigl[p_iV_A(T_A^s(s))
 +(1-p_i)V_A(T_A^f(s))\bigr]\Bigr\}\\
&+D_B(s)\gamma\Bigl[
p_{\pi_B(s)}V_A(T_B^s(s))
 +(1-p_{\pi_B(s)})V_A(T_B^f(s))\Bigr].
\end{aligned}
\end{equation}
With $\pi_A$ fixed, seller $B$'s Bellman equation is symmetric.

\paragraph{Equilibrium.}
A pair $(\pi_A^*,\pi_B^*)$ is a Markov perfect equilibrium if, at every
state $s$, $\pi_A^*(s)$ solves the maximization given $\pi_B^*$, and
$\pi_B^*(s)$ solves the corresponding maximization given $\pi_A^*$.
Equivalently,
\begin{equation}
\boxed{\pi_A^*\in\BR_A(\pi_B^*),\qquad
       \pi_B^*\in\BR_B(\pi_A^*).}
\end{equation}

\paragraph{Product 2.}
Conditional on being selected, we obtain the same threshold. Define
$\Delta c=c_2-c_1$ and $\Delta p=p_2-p_1$. For seller $A$,
\begin{equation}\label{eq:threshold}
\boxed{\text{Product 2 is optimal}
\quad\Longleftrightarrow\quad
\gamma\bigl[V_A(T_A^s(s))-V_A(T_A^f(s))\bigr]
\geq\frac{\Delta c}{\Delta p}.}
\end{equation}
At equality, both products are optimal.

\section{Initial intuition and computational concerns}

\paragraph{Initial intuition.}
Extinction should continue to hold because the customer will eventually
have accumulated so much information that her beliefs will change very
little. Thus, both sellers will eventually stop using product~2.

\paragraph{Problems and possible approaches.}
\begin{itemize}
\item In the base model, the state space is infinite (the state keeps
growing), which makes implementation computationally demanding.
\item Perhaps it would be better to work with forgetting.
\item If so, we could take advantage of the boundary.
\end{itemize}

\section{Does extinction occur?}

Not always, or more precisely, not at every state. Why? There are cases
in which the difference in demand does not tend to zero as $n$ grows.

\paragraph{Example.}
Consider the sequence of states
\begin{equation}
s_n=(0,n,0,n),\qquad n\in\mathbb N_0.
\end{equation}
In particular, the customer's draw for seller $B$ satisfies
$\widetilde\theta_B\sim\Beta(1,n+1)$. Calculate seller $A$'s demand
after each possible outcome:
\begin{align}
\text{Success:}\quad
D_A(T_A^s(s_n))
&=\Pb\bigl(\Beta(2,n+1)>\Beta(1,n+1)\bigr)
=\frac{3n+4}{4n+6},\label{eq:success-demand}\\[4pt]
\text{Failure:}\quad
D_A(T_A^f(s_n))
&=\Pb\bigl(\Beta(1,n+2)>\Beta(1,n+1)\bigr)
=\frac{n+1}{2n+3}.\label{eq:failure-demand}
\end{align}
The Beta notation inside these probabilities represents independent
random variables with the indicated distributions. Therefore,
\begin{equation}
\begin{aligned}
D_A(T_A^s(s_n))-D_A(T_A^f(s_n))
&=\frac{3n+4}{4n+6}-\frac{n+1}{2n+3}\\
&=\frac{n+2}{4n+6}
=\frac{1+2/n}{4+6/n}
\longrightarrow\frac14,
\end{aligned}
\end{equation}

In this case, because the posteriors concentrate around zero, seller $A$
may still benefit from using product~2. We cannot guarantee that the
difference in demand is small at every state simply because the total
number of observations is large.

\paragraph{Note.}
The calculation above concerns the demand difference which is the tool that we used in the base case to proof the extinction result. 
\eqref{eq:threshold}.

The thing with this example is that the probability of that state occurring is very small, and that is an important point. Maybe we need to study individual trajectories to understand what happens?

\paragraph{Computational problem.}
Because the state grows without bound, implementing the full problem is
infeasible or computationally very demanding. We would need to truncate
the problem.

\paragraph{What could we prove.}
Existence of equilibrium? 

\end{document}
